%%%%%%%%%%%%%%%%%
% This is an example CV created using altacv.cls (v1.1.5, 1 December 2018) written by
% LianTze Lim (liantze@gmail.com), based on the
% Cv created by BusinessInsider at http://www.businessinsider.my/a-sample-resume-for-marissa-mayer-2016-7/?r=US&IR=T
%
%% It may be distributed and/or modified under the
%% conditions of the LaTeX Project Public License, either version 1.3
%% of this license or (at your option) any later version.
%% The latest version of this license is in
%%    http://www.latex-project.org/lppl.txt
%% and version 1.3 or later is part of all distributions of LaTeX
%% version 2003/12/01 or later.
%%%%%%%%%%%%%%%%

%% If you are using \orcid or academicons
%% icons, make sure you have the academicons
%% option here, and compile with XeLaTeX
%% or LuaLaTeX.
% \documentclass[10pt,a4paper,academicons]{altacv}

%% Use the "normalphoto" option if you want a normal photo instead of cropped to a circle
% \documentclass[10pt,a4paper,normalphoto]{altacv}

\documentclass[10pt,a4paper,ragged2e]{altacv}

%% AltaCV uses the fontawesome and academicon fonts
%% and packages.
%% See texdoc.net/pkg/fontawecome and http://texdoc.net/pkg/academicons for full list of symbols. You MUST compile with XeLaTeX or LuaLaTeX if you want to use academicons.

% Change the page layout if you need to
\geometry{left=2cm,right=10cm,marginparwidth=6.8cm,marginparsep=1.2cm,top=1.25cm,bottom=1.25cm}

% Change the font if you want to, depending on whether
% you're using pdflatex or xelatex/lualatex
\ifxetexorluatex
  % If using xelatex or lualatex:
  \setmainfont{Carlito}
\else
  % If using pdflatex:
  \usepackage[utf8]{inputenc}
  \usepackage[T1]{fontenc}
  \usepackage[default]{lato}
\fi

% Change the colours if you want to
\definecolor{VividPurple}{HTML}{000000}
\definecolor{SlateGrey}{HTML}{2E2E2E}
\definecolor{LightGrey}{HTML}{2E2E2E}
\colorlet{heading}{VividPurple}
\colorlet{accent}{VividPurple}
\colorlet{emphasis}{SlateGrey}
\colorlet{body}{LightGrey}

% Change the bullets for itemize and rating marker
% for \cvskill if you want to
\renewcommand{\itemmarker}{{\small\textbullet}}
\renewcommand{\ratingmarker}{\faCircle}

%% Selected publications are maintained manually in page1sidebar.tex.
\usepackage[hidelinks]{hyperref}
\hypersetup{
    colorlinks=true,
    linkcolor=blue,
    filecolor=magenta,
    urlcolor=blue,
}

\begin{document}
\name{Skanda Vaidyanath}
\tagline{Founding AI Researcher, Yutori}
% Cropped to square from https://en.wikipedia.org/wiki/Marissa_Mayer#/media/File:Marissa_Mayer_May_2014_(cropped).jpg, CC-BY 2.0
%\photo{3.3cm}{profile.jpg}
\personalinfo{%
  % Not all of these are required!
  % You can add your own with \printinfo{symbol}{detail}
  \email{\href{mailto:svaidyan@stanford.edu}{svaidyan@stanford.edu}}
%   \phone{000-00-0000}
%  \mailaddress{Address, Street, 00000 County}
%  \location{Mumbai, India}
  \homepage{\href{https://skandavaidyanath.github.io/}{skandavaidyanath.github.io}}
%  \twitter{@marissamayer}
  \linkedin{\href{https://www.linkedin.com/in/skanda-vaidyanath}{skanda-vaidyanath}}
   \github{\href{https://github.com/skandavaidyanath}{skandavaidyanath}}
   \scholar{\href{https://scholar.google.com/citations?user=CS9tprgAAAAJ&hl=en&oi=ao}{Skanda Vaidyanath}}
   % I'm just making this up though.
%   \orcid{orcid.org/0000-0000-0000-0000} % Obviously making this up too. If you want to use this field (and also other academicons symbols), add "academicons" option to \documentclass{altacv}
}

%% Make the header extend all the way to the right, if you want.
\begin{fullwidth}
\makecvheader
\end{fullwidth}

%% Depending on your tastes, you may want to make fonts of itemize environments slightly smaller
\AtBeginEnvironment{itemize}{\small}

%% Provide the file name containing the sidebar contents as an optional parameter to \cvsection.
%% You can always just use \marginpar{...} if you do
%% not need to align the top of the contents to any
%% \cvsection title in the "main" bar.

\cvsection[page1sidebar]{Education}
\cvevent{M.S. Computer Science A.I. Track}{Stanford University}{Sep 2021 -- Jun 2023}{}
\begin{itemize}
    \item \textbf{CGPA: 4.06/4}
\end{itemize}
\cvevent{B.E. (Hons.) Computer Science with a Minor in Data Science}{BITS Pilani, Hyderabad Campus}{Jun 2016 -- May 2020}{}
\begin{itemize}
\item \textbf{Class valedictorian.}
%\smallskip
\item \textbf{CGPA: 9.93/10, Major CGPA (only CS courses): 10.00/10}
\end{itemize}



\cvsection{Recent Experience}

\cvevent{Founding AI Researcher}{Yutori}{Feb 2025 -- Present}{San Francisco, California}
\begin{itemize}
    \item Post-trained the \href{https://yutori.com/blog/introducing-n1}{Navigator n1}, \href{https://yutori.com/blog/introducing-n1-5}{n1.5}, and \href{https://yutori.com/blog/introducing-n2}{n2} model series using reinforcement learning and self-distillation (GRPO variants and \href{https://arxiv.org/abs/2601.20802}{SDPO}).
    \item \href{https://yutori.com/blog/introducing-n2}{Navigator n2} is a 27B computer-use model achieving \textbf{SOTA} results on four of five evaluated benchmarks.
    \item Built the agentic harness and backend for \href{https://yutori.com/scouts}{Scouts} and \href{https://yutori.com/delegate}{Delegate}.
\end{itemize}

\divider

\cvevent{Staff Research Scientist}{Riot Games AI Accelerator}{Jul 2023 -- Feb 2025}{Redwood City, California}
\begin{itemize}
    \item Trained self-play policies for two-player zero-sum partially observable stochastic games, achieving professional-level performance.
    \item Developed game-theoretic extensions to \href{https://arxiv.org/pdf/1912.00167}{asynchronous PPO} based on \href{https://arxiv.org/pdf/2206.15378}{R-NaD} and \href{https://arxiv.org/pdf/2206.05825}{magnetic mirror descent (MMD)}.
    \item Scaled self-play using large-scale distributed reinforcement learning.
\end{itemize}

\divider

\cvevent{Research Engineer Intern}{Google DeepMind}{Jun 2022 -- Sep 2022}{Mountain View, California}
\begin{itemize}
    \item Co-developed \href{https://deepmind-pushworld.github.io/play/}{PushWorld}, a benchmark of 200+ puzzles for manipulation planning with movable obstacles and tools \href{https://arxiv.org/abs/2301.10289}{[paper]}.
    \item Evaluated classical planners and model-free RL agents; introduced a planning heuristic 40$\times$ faster than the strongest baseline.
\end{itemize}

\divider

\cvevent{Research Intern}{Microsoft Research}{Dec 2020 -- Jul 2021}{Bangalore, India}
\begin{itemize}
    \item Built \href{https://www.microsoft.com/en-us/research/project/project-jigsaw/}{Jigsaw}, a multi-modal system that synthesizes Pandas code from natural language and I/O examples using GPT-3 and Codex \href{https://arxiv.org/html/2112.02969v1}{[paper]}.
    \item Developed program-analysis and synthesis post-processing to repair reference, argument, and semantic errors and learn from user feedback.
\end{itemize}

% \cvevent{Research Intern}{Max Planck Institute for Informatics}{Aug 2019 -- May 2020}{Saarbrucken, Germany}
% \begin{itemize}
%     \item Worked on fine-tuning LLMs with RL for information retrieval applications.
% \end{itemize}

% \cvevent{Research Intern}{Max Planck Institute for Informatics}{Aug 2019 -- May 2020}{Saarbrucken, Germany}
% \begin{itemize}
%     %\item Used deep reinforcement learning to recommend travel destinations to users based on forum posts and user reviews.
%     %\item Used deep reinforcement learning to improve document retrieval from a corpus via document expansion.
%     \item Worked on projects at the intersection of deep RL and information retrieval
%     \item \textbf{Keywords:} Information retrieval, natural language processing, reinforcement learning
% \end{itemize}

% \cvevent{Research Intern (IUSSTF-Viterbi programme)}{USC Institute for Creative Technologies}{May 2019 -- July 2019}{Los Angeles, USA}
% \begin{itemize}
%     %\item Developed a policy to control a swarm of drones that attempts to negotiate with and save civilians from a forest fire and achieved a success rate of over 95\% on the task.
%     \item Worked on RL algorithms to control a swarm of drones that attempts to negotiate with and save civilians from a forest fire.
%     \item \textbf{Keywords:} Reinforcement learning
% \end{itemize}

%\divider

% \divider
%\cvskill{German}{3}


%\divider
% \divider

% \cvevent{Product Engineer}{Google}{23 June 1999 -- 2001}{Palo Alto, CA}

% \begin{itemize}
% \item Joined the company as employe \#20 and female employee \#1
% \item Developed targeted advertisement in order to use user's search queries and show them related ads
% \end{itemize}

%\cvsection{A Day of My Life}

% Adapted from @Jake's answer from http://tex.stackexchange.com/a/82729/226
% \wheelchart{outer radius}{inner radius}{
% comma-separated list of value/text width/color/detail}
% Some ad-hoc tweaking to adjust the labels so that they don't overlap
% \wheelchart{1.5cm}{0.5cm}{%
%   10/10em/accent!30/Sleeping \& dreaming about work,
%   25/9em/accent!60/Public resolving issues with Yahoo!\ investors,
%   5/13em/accent!10/\footnotesize\\[1ex]New York \& San Francisco Ballet Jawbone board member,
%   20/15em/accent!40/Spending time with family,
%   5/8em/accent!20/\footnotesize Business development for Yahoo!\ after the Verizon acquisition,
%   30/9em/accent/Showing Yahoo!\ employees that their work has meaning,
%   5/8em/accent!20/Baking cupcakes
% }

\clearpage

% \cvsection[page2sidebar]{Publications}

% \printbibliography[heading=pubtype,title={\printinfo{\faBook}{Books}},type=book]

% \divider

% \printbibliography[heading=pubtype,title={\printinfo{\faFileTextO}{Journal Articles}}, type=article]

% \divider

% \printbibliography[heading=pubtype,title={\printinfo{\faGroup}{Conference Proceedings}},type=inproceedings]

% %% If the NEXT page doesn't start with a \cvsection but you'd
% %% still like to add a sidebar, then use this command on THIS
% %% page to add it. The optional argument lets you pull up the
% %% sidebar a bit so that it looks aligned with the top of the
% %% main column.
% % \addnextpagesidebar[-1ex]{page3sidebar}


\end{document}
